% !TEX program = pdflatex
% ------------------------------------------------------------
% MATH 431 Homework Template (plug-and-chug friendly)
%
% Goals:
% - Same clean header style as your MATH 521 template.
% - Easy “Problem -> Solution” workflow.
% - Lightweight (no tcolorbox); uses mdframed for optional boxes.
% - Good for probability homework: sample space, probability measure, events.
%
% Quick use:
%   1) Edit the Metadata block.
%   2) For each exercise, use \problem{<section>.<number>} and paste the prompt.
%   3) Write your solution in the solution environment.
% ------------------------------------------------------------

\documentclass[12pt]{article}

% =========================
% 0) Metadata (EDIT ME)
% =========================
\newcommand{\CourseCode}{MATH 431}
\newcommand{\CourseTitle}{Introduction to Probability}
\newcommand{\CourseSection}{LEC 002} % change to 001 or 003 if needed
\newcommand{\Semester}{Spring 2026}
\newcommand{\Instructor}{Sarah Tammen} % LEC 001: Ander Aguirre | LEC 002: Sarah Tammen | LEC 003: Nicholas Derr
\newcommand{\MeetingInfo}{MWF 9:55--10:45, 6102 Sewell Social Sciences} % LEC 001 default; update if in 002/003
\newcommand{\HomeworkNumber}{3}
\newcommand{\YourName}{Alejandro De La Torre}
\newcommand{\DueDate}{February 13, 2026} % or e.g. February 2, 2026

% =========================
% 1) Core math tools
% =========================
\usepackage{amsmath, amssymb, amsthm}

% =========================
% 2) Lists and spacing
% =========================
\usepackage{enumitem}
\setlist[enumerate,1]{label=\textbf{(\alph*)}, leftmargin=2em, itemsep=0.25em}
\setlist[itemize]{leftmargin=1.5em, itemsep=0.25em}
\setlength{\parskip}{0.35em}
\setlength{\parindent}{0pt}

% =========================
% 3) Page geometry + header/footer
% =========================
\usepackage{geometry}
\geometry{margin=1in}

\usepackage{fancyhdr}
\pagestyle{fancy}
\fancyhf{}
\lhead{\CourseCode\ --- \CourseSection, \Semester}
\rhead{Homework \HomeworkNumber}
\cfoot{\thepage}
\renewcommand{\headrulewidth}{0.4pt}
\setlength{\headheight}{15pt}
\addtolength{\topmargin}{-2pt}

% =========================
% 4) Links (subtle)
% =========================
\usepackage[hidelinks]{hyperref}
\setcounter{tocdepth}{1}

% =========================
% 5) Figures (optional)
% =========================
\usepackage{graphicx}
\graphicspath{{images/}}

% =========================
% 6) Lightweight boxes (optional)
% =========================
\usepackage{mdframed}
\newmdenv[
  skipabove=0.6\baselineskip,
  skipbelow=0.6\baselineskip,
  linewidth=0.6pt,
  roundcorner=4pt,
  innerleftmargin=10pt,
  innerrightmargin=10pt,
  innertopmargin=8pt,
  innerbottommargin=8pt
]{boxnote}

% =========================
% 7) Probability / notation macros
% =========================
\newcommand{\R}{\mathbb{R}}
\newcommand{\N}{\mathbb{N}}
\newcommand{\Z}{\mathbb{Z}}

\newcommand{\OmegaSpace}{\Omega}
\newcommand{\FField}{\mathcal{F}}
\newcommand{\E}{\mathbb{E}}

\newcommand{\given}{\,\middle|\,}
\newcommand{\ind}{\mathbf{1}}

% Common “defined as” symbol
\newcommand{\defeq}{\vcentcolon=}

% =========================
% 8) Problem command (TOC + PDF bookmarks)
% Usage:
%   \problem{<label>}
%   \problem[Short title]{<label>}
% Example:
%   \problem{1.4}
%   \problem[Sample space]{1.4}
% =========================
\makeatletter
\newcommand{\problem}[2][]{%
  \def\prob@title{Problem #2\if\relax\detokenize{#1}\relax\else\ (\textit{#1})\fi}%
  \phantomsection%
  \pdfbookmark[1]{\prob@title}{prob#2}%
  \addcontentsline{toc}{section}{\prob@title}%
  \section*{\prob@title}%
  \vspace{-0.25em}\hrule\vspace{0.8em}%
}
\makeatother

% =========================
% 9) Solution environment
% =========================
\newenvironment{solution}{%
  \noindent\textbf{Solution.}\ \ignorespaces
}{\par}

% Optional scratch / planning block
\newenvironment{scratch}{%
  \begin{boxnote}\textbf{Scratch / Setup.}\ \small
}{\end{boxnote}}

% Final answer command: expects math only (no $$ or \[ \])
\newcommand{\finalanswer}[1]{%
  \par\medskip
  \noindent\textbf{Final:}\ \fbox{$#1$}\par
} % AGAIN: MATH ONLY INSIDE THE BOX; if word answer, use \text{} inside or use \fbox{TEXT} or \framebox

% =========================
% 10) Title block
% =========================
\title{Homework \HomeworkNumber}
\author{\YourName \\
\CourseCode\ --- \CourseTitle \\
\CourseSection \\
Instructor: \Instructor}
\date{Due: \DueDate}

\begin{document}
\maketitle

% \section*{Collaboration / Resources}
% (Optional) Write “I worked alone” or list collaborators/resources.

\tableofcontents
\begingroup
\renewcommand{\thefootnote}{}
\footnotetext{Notation: We sometimes use the standard shorthand for intersections of events: $AB$ means $A\cap B$, $A^cB$ means $A^c\cap B$, and $ABC$ means $A\cap B\cap C$. This shorthand is used consistently throughout.}
\endgroup
\setcounter{footnote}{0}
\bigskip
\hrule
\bigskip
\newpage

% ============================================================
% Problems (duplicate blocks as needed)
% ============================================================

\problem{2.4}
We have two urns. The first urn contains two balls labeled 1 and 2. The second urn contains three balls labeled 3, 4 and 5. We choose one of the urns at random (with equal probability) and then sample one ball (uniformly at random) from the chosen urn. What is the probability that we picked the ball labeled 5?

\begin{solution}
\begin{scratch}
Let $U_1$ be the event we choose urn 1, $U_2$ the event we choose urn 2. Let $F$ be the event that we pick ball 5. Use the law of total probability over $U_1, U_2$.
\end{scratch}

We have $P(U_1)=P(U_2)=\tfrac12$. Also $P\!\left(F \given U_1\right)=0$ and $P\!\left(F \given U_2\right)=\tfrac13$. Therefore,
\[
P(F)=P\!\left(F \given U_1\right)P(U_1)+P\!\left(F \given U_2\right)P(U_2)
=0\cdot \tfrac12+\tfrac13\cdot \tfrac12=\tfrac16.
\]
\finalanswer{\tfrac16}
\end{solution}

\newpage

\problem{2.10}
I have a bag with 3 fair dice. One is 4-sided, one is 6-sided, and one is 12-sided. I reach into the bag, pick one die at random and roll it. The outcome of the roll is 4. What is the probability that I pulled out the 6-sided die?

\begin{solution}
\begin{scratch}
Let $D_4, D_6, D_{12}$ be the events that the 4-, 6-, or 12-sided die is chosen. Let $R$ be the event that the roll is 4. Use Bayes' rule.
\end{scratch}

Each die is chosen with probability $1/3$. Likelihoods: $P\!\left(R \given D_4\right)=\tfrac14$, $P\!\left(R \given D_6\right)=\tfrac16$, $P\!\left(R \given D_{12}\right)=\tfrac1{12}$. Hence
\[
P\!\left(D_6 \given R\right)=\frac{P\!\left(R \given D_6\right)P(D_6)}
{P\!\left(R \given D_4\right)P(D_4)+P\!\left(R \given D_6\right)P(D_6)+P\!\left(R \given D_{12}\right)P(D_{12})}
=\frac{\tfrac16\cdot \tfrac13}{\tfrac14\cdot \tfrac13+\tfrac16\cdot \tfrac13+\tfrac1{12}\cdot \tfrac13}
=\tfrac13.
\]
\finalanswer{\tfrac13}
\end{solution}

\newpage

\problem{2.14}
Let $A$ and $B$ be two disjoint events. Under what condition are they independent?

\begin{solution}
\begin{scratch}
Disjoint means $A\cap B=\varnothing$, so $P(AB)=0$. Independence requires $P(AB)=P(A)P(B)$.
\end{scratch}

Since $A$ and $B$ are disjoint, $P(AB)=0$. Independence requires $P(A)P(B)=0$, which happens iff at least one of $P(A)$ or $P(B)$ is zero. 
\finalanswer{P(A)=0\ \text{or}\ P(B)=0}
\end{solution}

\newpage

\problem{2.16}
We flip a fair coin three times. For $i=1,2,3$, let $A_i$ be the event that among the first $i$ coin flips we have an odd number of heads. Check whether the events $A_1, A_2, A_3$ are independent or not.

\begin{solution}
\begin{scratch}
Let $\Omega=\{H,T\}^3$ with $|\Omega|=8$. Compute $P(A_i)$ and all intersections of the $A_i$.
\end{scratch}

We have $P(A_1)=P(\text{first flip is }H)=\tfrac12$, $P(A_2)=\tfrac12$ (odd heads among first two), and $P(A_3)=\tfrac12$ (odd heads among three).

Intersections: 
\[
P(A_1A_2)=P(HT)=\tfrac14,\quad
P(A_1A_3)=P(\{HHH,HTT\})=\tfrac14,
\]
\[
P(A_2A_3)=P(\{HTT,THT\})=\tfrac14,\quad
P(A_1A_2A_3)=P(HTT)=\tfrac18.
\]

Each pair satisfies $P(A_iA_j)=P(A_i)P(A_j)=\tfrac14$, and the triple satisfies
$P(A_1A_2A_3)=P(A_1)P(A_2)P(A_3)=\tfrac18$.
Therefore the events are mutually independent.
\finalanswer{\text{$A_1, A_2, A_3$ are mutually independent}}
\end{solution}

\newpage

\problem{2.18}
We choose a number from the set $\{10, 11, 12, \ldots, 99\}$ uniformly at random.
\begin{enumerate}[label=(\alph*)]
\item Let $X$ be the first digit and $Y$ the second digit of the chosen number. Show that $X$ and $Y$ are independent random variables.
\item Let $X$ be the first digit of the chosen number and $Z$ the sum of the two digits. Show that $X$ and $Z$ are not independent.
\end{enumerate}

\begin{solution}
\begin{scratch}
There are $90$ equally likely numbers. For (a), compute $P(X=x)$, $P(Y=y)$, and $P(X=x,Y=y)$. For (b), give a specific $(x,z)$ where factorization fails.
\end{scratch}

(a) For $x\in\{1,\ldots,9\}$ and $y\in\{0,\ldots,9\}$, there is exactly one number with first digit $x$ and second digit $y$. Thus
\[
P(X=x,Y=y)=\tfrac1{90}.
\]
Also $P(X=x)=\tfrac{10}{90}=\tfrac19$ and $P(Y=y)=\tfrac{9}{90}=\tfrac1{10}$. Hence
\[
P(X=x,Y=y)=\tfrac1{90}=\tfrac19\cdot\tfrac1{10}=P(X=x)P(Y=y),
\]
so $X$ and $Y$ are independent.

(b) Take $x=1$ and $z=1$. Then $Z=1$ occurs only for the number $10$, so $P(Z=1)=\tfrac1{90}$. We have
\[
P(X=1,Z=1)=\tfrac1{90},\quad P(X=1)P(Z=1)=\tfrac19\cdot\tfrac1{90}=\tfrac1{810},
\]
which are not equal. Thus $X$ and $Z$ are not independent.
\finalanswer{\text{$X$ and $Y$ are independent; $X$ and $Z$ are not independent}}
\end{solution}

\newpage

\problem{2.30}
Assume that $1/3$ of all twins are identical twins. You learn that Miranda is expecting twins, but you have no other information.
\begin{enumerate}[label=(\alph*)]
\item Find the probability that Miranda will have two girls.
\item You learn that Miranda gave birth to two girls. What is the probability that the girls are identical twins?
\end{enumerate}
Explain any assumptions you make.

\begin{solution}
\begin{scratch}
Assume identical twins have the same sex, with boy/girl equally likely; fraternal twins have independent sexes with boy/girl equally likely. Let $I$ be identical, $F$ fraternal, and $G$ be the event of two girls.
\end{scratch}

We have $P(I)=\tfrac13$ and $P(F)=\tfrac23$. Under the assumptions, $P\!\left(G \given I\right)=\tfrac12$ and $P\!\left(G \given F\right)=\tfrac14$.

(a) By total probability,
\[
P(G)=P\!\left(G \given I\right)P(I)+P\!\left(G \given F\right)P(F)
=\tfrac12\cdot\tfrac13+\tfrac14\cdot\tfrac23=\tfrac13.
\]
\finalanswer{P(G)=\tfrac13}

(b) By Bayes' rule,
\[
P\!\left(I \given G\right)=\frac{P\!\left(G \given I\right)P(I)}{P(G)}
=\frac{\tfrac12\cdot\tfrac13}{\tfrac13}=\tfrac12.
\]
\finalanswer{P\!\left(I \given G\right)=\tfrac12}
\end{solution}

\newpage

\problem{2.34}
You play the following game against your friend. You have two urns and three balls. One of the balls is marked. You get to place the balls in the two urns any way you please, including leaving one urn empty. Your friend will choose one urn at random and then draw a ball from that urn. (If he chose an empty urn, there is no ball.) His goal is to draw the marked ball.
\begin{enumerate}[label=(\alph*)]
\item How would you arrange the balls in the urns to \emph{minimize} his chances of drawing the marked ball?
\item How would your friend arrange the balls in the urns to \emph{maximize} his chances of drawing the marked ball?
\item Repeat (a) and (b) for the case of $n$ balls with one marked ball.
\end{enumerate}

\begin{solution}
\begin{scratch}
Let $M$ be the event the marked ball is drawn. If the marked ball is in an urn with $k$ unmarked balls (so $k+1$ total), and the other urn has the remaining balls (possibly empty), then $P(M)=\tfrac12\cdot \tfrac1{k+1}$.
\end{scratch}

If the marked ball is in an urn with $k$ unmarked balls, then the friend chooses that urn with probability $\tfrac12$ and, conditional on that, draws the marked ball with probability $\tfrac1{k+1}$. Hence
\[
P(M)=\tfrac12\cdot \tfrac1{k+1}.
\]
(a) To minimize $P(M)$ we maximize $k$. With three balls total, the best is to put all three balls in one urn and leave the other empty ($k=2$), giving $P(M)=\tfrac12\cdot \tfrac13=\tfrac16$.

(b) To maximize $P(M)$ we minimize $k$, so put the marked ball alone in one urn and the two unmarked balls in the other ($k=0$), giving $P(M)=\tfrac12$.

(c) With $n$ balls total, the same calculation gives $P(M)=\tfrac12\cdot \tfrac1{k+1}$ when the marked ball is grouped with $k$ unmarked balls. Thus
\[
\text{minimize: }k=n-1\Rightarrow P(M)=\tfrac1{2n},\qquad
\text{maximize: }k=0\Rightarrow P(M)=\tfrac12.
\]
\finalanswer{\text{minimize } P(M)=\tfrac1{2n}\ \text{(all balls together); maximize } P(M)=\tfrac12\ \text{(marked ball alone)}}
\end{solution}

\newpage

\problem{2.52}
Suppose that $P(A)=0.3$, $P(B)=0.2$, and $P(C)=0.1$. Further, $P(A \cup B)=0.44$, $P(A^cC)=0.07$, $P(BC)=0.02$, and $P(A \cup B \cup C)=0.496$. Decide whether $A$, $B$, and $C$ are mutually independent.

\begin{solution}
\begin{scratch}
Compute $P(AB)$ from $P(A\cup B)$, compute $P(AC)$ from $P(A^cC)$, and compare pairwise products. Then use inclusion--exclusion to compute $P(ABC)$ and compare to $P(A)P(B)P(C)$.
\end{scratch}

First,
\[
P(AB)=P(A)+P(B)-P(A\cup B)=0.3+0.2-0.44=0.06.
\]
Also $P(A^cC)=0.07$, so $P(AC)=P(C)-P(A^cC)=0.1-0.07=0.03$.
We are given $P(BC)=0.02$.

These match pairwise products: $0.3\cdot 0.2=0.06$, $0.3\cdot 0.1=0.03$, and $0.2\cdot 0.1=0.02$.

Now use inclusion--exclusion:
\[
P(A\cup B\cup C)=P(A)+P(B)+P(C)-P(AB)-P(AC)-P(BC)+P(ABC).
\]
Thus
\[
P(ABC)=0.496-(0.3+0.2+0.1)+(0.06+0.03+0.02)=0.006.
\]
Since $P(A)P(B)P(C)=0.3\cdot 0.2\cdot 0.1=0.006$, the triple also factorizes. Hence the events are mutually independent.
\finalanswer{\text{$A,B,C$ are mutually independent}}
\end{solution}

\newpage

\problem{2.54}
Let $A$ and $B$ be events with these properties: $0 < P(B) < 1$ and $P(A|B)=P(A|B^c)=1/3$.
\begin{enumerate}[label=(\alph*)]
\item Is it possible to calculate $P(A)$ from this information? Either declare that it is not possible, or find the value of $P(A)$.
\item Are $A$ and $B$ independent, not independent, or is it impossible to determine?
\end{enumerate}

\begin{solution}
\begin{scratch}
Use the law of total probability for $P(A)$, then compare $P\!\left(A \given B\right)$ to $P(A)$ for independence.
\end{scratch}

(a) By total probability,
\[
P(A)=P\!\left(A \given B\right)P(B)+P\!\left(A \given B^c\right)P(B^c)
=\tfrac13 P(B)+\tfrac13 P(B^c)=\tfrac13.
\]
\finalanswer{P(A)=\tfrac13}

(b) Since $P\!\left(A \given B\right)=\tfrac13=P(A)$, the events are independent.
\finalanswer{\text{$A$ and $B$ are independent}}
\end{solution}

\newpage

\end{document}
